\chapter{Befehlsreferenz}\label{kap:befehle}
\sloppy % lange Befehls-Aufforderungen in Schreibmaschinenschrift umbrechen
\index{Befehlsreferenz}

Dieser Anhang beschreibt jeden Befehl der Befehlszeile in alphabetischer Reihenfolge seines
Namens. Er ergänzt Kapitel~\ref{kap:modellieren}, das die Befehle mit Hintergrund und Beispielen
erklärt, um eine kompakte Nachschlagetabelle. Alle Aufforderungen (Prompts) und Meldungen sind
so wiedergegeben, wie sie im Programm erscheinen; Werte in Klammern wie \code{(6 sides)} zeigen
die Voreinstellung und ändern sich mit Ihrer Eingabe.

% ============================================================================
\section{Eingabesyntax der Befehlszeile}\label{sec:eingabesyntax}
\index{Eingabesyntax}

Was Sie in die Befehlszeile tippen und mit \keys{Enter} bestätigen, wird nach folgenden Regeln
gedeutet. Leerzeichen werden vorher entfernt, Groß- und Kleinschreibung spielt keine Rolle.
Tippen in einer Ansicht schreibt automatisch in die Befehlszeile (Ausnahmen: \keys{Z},
\keys{Leertaste} und \keys{H}, solange etwas ausgewählt ist).

\begin{referenz}{>{\raggedright\arraybackslash}p{0.25\linewidth}X}
\toprule
Eingabe & Bedeutung \\
\midrule\endhead
(leer) & laufenden Befehl beenden; läuft keiner, den zuletzt gestarteten Befehl wiederholen.
  Ein \maus{Rechtsklick} ohne Ziehen in einer Ansicht wirkt ebenso. \\
\code{line}, \code{l} \ldots & Beginnt die Eingabe mit einem Buchstaben, ist sie ein
  Befehlsname oder Alias. Ein laufender Befehl wird abgebrochen und der neue gestartet. \\
\code{12.5}, \code{-3}, \code{1e-3} & eine Zahl (siehe unten). \\
\code{x,y,z} & Punkt in absoluten Weltkoordinaten. \\
\code{x,y} & Punkt in der Konstruktionsebene der Ansicht, über der die Maus während des
  Befehls zuletzt war (vorher: der aktiven Ansicht): $\text{Ursprung} + x \cdot u + y \cdot v$. \\
\code{@dx,dy,dz} & Punkt relativ zum Bezugspunkt des Befehls, Versatz in Weltkoordinaten. \\
\code{@dx,dy} & Punkt relativ zum Bezugspunkt, Versatz in der Konstruktionsebene. \\
\code{d<a} & Polarkoordinaten in der Konstruktionsebene: Abstand $d$ vom Ursprung, Winkel $a$
  in Grad von der Achse $u$ zur Achse $v$. \\
\code{@d<a} & Polarkoordinaten relativ zum Bezugspunkt. \\
\bottomrule
\end{referenz}

\paragraph{Zahlenformat.} Dezimaltrennzeichen ist der Punkt (\code{2.5}, nicht \code{2,5}),
Koordinaten werden durch Kommas getrennt. Vorzeichen und Exponentenschreibweise
(\code{-1.5e2}) sind erlaubt. Ein Komma in einer einzelnen Zahl macht daraus eine
Koordinateneingabe.

\paragraph{Achsen der Konstruktionsebenen.}
\begin{referenz}{lllll}
\toprule
Ansicht & $u$ & $v$ & Normale $n$ & \code{3,4} ergibt \\
\midrule\endhead
Top, Perspective & $+x$ & $+y$ & $+z$ & $(3,4,0)$ \\
Front & $+x$ & $+z$ & $-y$ & $(3,0,4)$ \\
Left & $-y$ & $+z$ & $-x$ & $(0,-3,4)$ \\
\bottomrule
\end{referenz}

\paragraph{Zahlen.} Eine Zahl wird verwendet für das, wonach die Aufforderung fragt (Radius,
Höhe, Anzahl, Winkel, Faktor, Abstand). Kann der aktuelle Schritt keine Zahl verwenden, aber
einen Punkt, so ist die Zahl eine \emph{Länge}: Der Punkt entsteht in diesem Abstand vom
Bezugspunkt in Richtung des Cursors (bei eingeschaltetem Ortho in der korrigierten Richtung).
Die Länge muss positiv sein.

\paragraph{Bezugspunkt.} Relative Eingaben, Längen und Ortho beziehen sich auf den Bezugspunkt;
welcher Punkt das ist, steht bei jedem Befehl unten. Ohne Bezugspunkt wird \code{@\ldots} mit
\code{relative input needs a previous point} abgelehnt.

\paragraph{Weitere Eingaben während eines Befehls.}
\begin{referenz}{>{\raggedright\arraybackslash}p{0.25\linewidth}X}
\toprule
Eingabe & Wirkung \\
\midrule\endhead
\maus{Linksklick} & Punkt setzen (mit Objektfang). \\
\maus{Rechtsklick} ohne Ziehen, \keys{Enter} & Befehl beenden (ohne laufenden Befehl: den
  letzten wiederholen). Rechts \emph{ziehen} dreht bzw.\ verschiebt nur. \\
\keys{Esc} & Befehl abbrechen. \\
\keys{F8} & Ortho ein/aus. \\
\keys{F3} & Objektfang ein/aus. \\
\keys{Tab}, \keys{Shift+Tab} & nächster/vorheriger Fangkandidat. \\
\keys{Alt} & Fang unterdrücken; \keys{Alt}+\maus{Klick}: freier Punkt. \\
\bottomrule
\end{referenz}

\paragraph{Meldungen der Befehlszeile.}
\begin{referenz}{>{\raggedright\arraybackslash}p{0.38\linewidth}X}
\toprule
Meldung & Bedeutung \\
\midrule\endhead
\code{Unknown command '<name>'} & kein Befehl und kein Alias dieses Namens. \\
\code{<Titel>: <Grund>} & Eingabe abgelehnt oder Ergebnis nicht berechenbar, mit dem Grund
  (z.\,B.\ \code{Fillet: radius 5 too large for these curves}); ein wartender Befehl wartet
  weiter. \\
\code{Input not accepted -- <prompt>} & Eingabe in diesem Schritt ungültig, ohne eigenen Grund
  des Befehls; der Befehl wartet weiter. \\
\code{No command is waiting for a point} & Koordinaten ohne laufenden Befehl. \\
\code{A number is not expected here} & Zahl, die weder als Wert noch als Länge verwendbar
  ist. \\
\code{relative input needs a previous point} & \code{@} ohne Bezugspunkt. \\
\code{polar input: distance<angle} & fehlerhafte Polar-Eingabe. \\
\code{expected x,y or x,y,z} & falsche Anzahl von Koordinaten. \\
\code{'<text>' is not a number} & ein Koordinatenwert ist keine Zahl. \\
\code{<Titel> cancelled} & Befehl ohne Ergebnis beendet oder abgebrochen. \\
\code{<Titel>: select objects first} & Auswahlbefehl ohne Auswahl gestartet. \\
\code{Boolean <\ldots> failed: <Grund>} & eine Boolesche Operation ist gescheitert. \\
\bottomrule
\end{referenz}

\paragraph{Auswahlbefehle.} Befehle mit dem Eintrag \enquote{Auswahl} in der Übersicht arbeiten
auf den \emph{vor} dem Start ausgewählten ganzen Objekten, in Auswahlreihenfolge; einzeln
gewählte Flächen, Kanten und Ecken werden übergangen, ebenso Objekte unpassenden Typs.
\enquote{Sofort} bedeutet: Der Befehl braucht keine weitere Eingabe und ist mit dem Start
fertig.

\paragraph{Rückgängig.} Jedes Ergebnis ist ein Rückgängig-Schritt mit dem Titel des Befehls;
bei Point und Copy ist jeder Punkt bzw.\ jede Kopie ein eigener Schritt.

% ============================================================================
\section{Übersicht}\label{sec:befehle-uebersicht}

\begin{referenz}{ll>{\raggedright\arraybackslash}X>{\raggedright\arraybackslash}p{0.13\linewidth}}
\toprule
Name & Aliasse & Werkzeug (Gruppe, Knopf) & Eingabe \\
\midrule\endhead
\code{angle} & \code{ang} & Analyze, Angle (Ang) & Punkte \\
\code{arc} & \code{a} & Curves, Arc (Arc) & Punkte \\
\code{area} & & Analyze, Area (Are) & Auswahl, sofort \\
\code{array\_grid} & & Transform, Grid array (ArG) & Auswahl \\
\code{array\_linear} & \code{array} & Transform, Linear array (ArL) & Auswahl \\
\code{array\_polar} & \code{arraypolar} & Transform, Polar array (ArP) & Auswahl \\
\code{box} & \code{b} & Solids, Box (Box) & Punkte \\
\code{cap} & & Surfaces, Cap planar holes (Cap) & Auswahl, sofort \\
\code{circle} & \code{c}, \code{ci} & Curves, Circle (centre, radius) (Ci) & Punkte \\
\code{circle3p} & \code{c3} & Curves, Circle (3 points) (C3) & Punkte \\
\code{cone} & & Solids, Cone (Con) & Punkte \\
\code{coons} & & Surfaces, Coons patch (Cns) & Auswahl, sofort \\
\code{copy} & \code{co}, \code{cp} & Transform, Copy (Cp) & Auswahl \\
\code{curvature} & \code{curv} & Analyze, Curvature (Crv) & Auswahl \\
\code{cylinder} & \code{cyl} & Solids, Cylinder (Cyl) & Punkte \\
\code{difference} & \code{bd} & Solids, Boolean difference ($-$) & Auswahl, sofort \\
\code{distance} & \code{dist}, \code{di} & Analyze, Distance (Dst) & Punkte \\
\code{explode} & \code{x} & Edit, Explode (Exp) & Auswahl, sofort \\
\code{extrude} & \code{ext} & Surfaces, Extrude (Ext) & Auswahl \\
\code{fillet\_curve} & \code{fillet} & Curves, Fillet curves (Fil) & Auswahl \\
\code{fillet\_surface} & \code{filletsrf} & Surfaces, Fillet surfaces (FiS) & Auswahl \\
\code{gordon} & \code{network} & Surfaces, Gordon surface (Gor) & Auswahl, sofort \\
\code{helix} & \code{h}, \code{he} & Curves, Helix (Hx) & Punkte \\
\code{intersection} & \code{bi} & Solids, Boolean intersection ($\cap$) & Auswahl, sofort \\
\code{intersection\_check} & \code{intersect} & Analyze, Intersection check (Chk) & Auswahl,
  sofort \\
\code{join} & \code{j} & Edit, Join (Jn) & Auswahl, sofort \\
\code{line} & \code{l}, \code{li} & Draw, Line (Ln) & Punkte \\
\code{loft} & & Surfaces, Loft (Lft) & Auswahl, sofort \\
\code{mass\_properties} & \code{volume} & Analyze, Volume / mass properties (Vol) & Auswahl,
  sofort \\
\code{mirror} & \code{mi} & Transform, Mirror (Mir) & Auswahl \\
\code{move} & \code{m} & Transform, Move (Mv) & Auswahl \\
\code{offset\_curve} & \code{offset} & Curves, Offset curve (Off) & Auswahl \\
\code{offset\_surface} & \code{offsetsrf} & Surfaces, Offset surface (OfS) & Auswahl \\
\code{plane} & & Surfaces, Plane (Pln) & Punkte \\
\code{point} & \code{pt}, \code{po} & Draw, Point (Pt) & Punkte \\
\code{polygon} & \code{pol} & Draw, Polygon (Pol) & Punkte \\
\code{polyline} & \code{pl} & Draw, Polyline (PL) & Punkte \\
\code{rectangle} & \code{rec}, \code{rect} & Draw, Rectangle (Rec) & Punkte \\
\code{revolve} & \code{rev} & Surfaces, Revolve (Rev) & Auswahl \\
\code{rotate} & \code{ro} & Transform, Rotate (Rot) & Auswahl \\
\code{scale} & \code{sc} & Transform, Scale (Scl) & Auswahl \\
\code{sphere} & \code{sph} & Solids, Sphere (Sph) & Punkte \\
\code{spline} & \code{sp} & Curves, Interpolated spline (Sp) & Punkte \\
\code{split} & & Edit, Split (Spl) & Auswahl \\
\code{sweep1} & & Surfaces, Sweep 1 rail (Sw1) & Auswahl, sofort \\
\code{sweep2} & & Surfaces, Sweep 2 rails (Sw2) & Auswahl, sofort \\
\code{torus} & \code{tor} & Solids, Torus (Tor) & Punkte \\
\code{trim} & \code{tr} & Edit, Trim (Trm) & Auswahl \\
\code{union} & \code{bu} & Solids, Boolean union ($\cup$) & Auswahl, sofort \\
\midrule
-- & & Edit, Delete (Del); Taste \keys{Entf} & Auswahl, sofort \\
\bottomrule
\end{referenz}

In Klammern steht die Kurzbeschriftung des Knopfs in der Werkzeugspalte. Die Befehle haben
keine eigenen Tastenkürzel; der Tooltip jedes Werkzeugs nennt stattdessen, was Sie tippen
(z.\,B.\ \enquote{Type: l + Enter}). Das Werkzeug \gui{Delete} ist kein Befehl der
Befehlszeile; es entspricht \menu{Edit,Delete} und hat als einziges ein Tastenkürzel
(\keys{Entf}).

% ============================================================================
\section{Befehle von A bis Z}\label{sec:befehle-az}

% Darstellung eines Eintrags
\newcommand{\befehlseintrag}[1]{\par\begingroup\setlength{\parfillskip}{0pt}\noindent\begin{tabularx}{\linewidth}{@{}>{\bfseries}p{0.2\linewidth}X@{}}#1\end{tabularx}\par\endgroup}

\subsection{angle -- Angle}\label{sec:befehl-angle}
\befehlseintrag{%
Name & \cmd{angle}; Alias \code{ang} \\
Werkzeug & Analyze, \gui{Angle} \\
Eingabe & 1.~\code{Angle: vertex}\newline
  2.~\code{Angle: point on the first leg}\newline
  3.~\code{Angle: point on the second leg} \\
Bezugspunkt & der Scheitel \\
Ergebnis & Meldung \code{Angle <w>}\si{\degree} -- räumlicher Winkel zwischen den Schenkeln, 0 bis
  \SI{180}{\degree}. Es wird nichts erzeugt. \\
Hinweise & Ein Schenkelpunkt, der mit dem Scheitel zusammenfällt, wird abgelehnt. \\}

\subsection{arc -- Arc}\label{sec:befehl-arc}
\befehlseintrag{%
Name & \cmd{arc}; Alias \code{a} \\
Werkzeug & Curves, \gui{Arc} \\
Eingabe & 1.~\code{Arc: start point}\newline
  2.~\code{Arc: end point}\newline
  3.~\code{Arc: point on the arc} \\
Bezugspunkt & der zuletzt gesetzte Punkt \\
Ergebnis & Kreisbogen vom Start- über den dritten zum Endpunkt, in beliebiger Raumlage. \\
Hinweise & Doppelte Punkte und drei Punkte auf einer Geraden werden abgelehnt. \\}

\subsection{area -- Area / length}\label{sec:befehl-area}
\befehlseintrag{%
Name & \cmd{area} \\
Werkzeug & Analyze, \gui{Area} \\
Auswahl & Kurven, Flächen, Körper; sofort \\
Ergebnis & Meldung \code{Area <A>   Length <L>}: Summe der Flächeninhalte aller Flächen
  (bei Körpern aller Teilflächen) und Summe der Kurvenlängen; nur vorhandene Teile werden
  angezeigt. \\
Fehler & \code{Area / length: no curves or surfaces selected} \\}

\subsection{array\_grid -- Grid array}\label{sec:befehl-array-grid}
\befehlseintrag{%
Name & \cmd{array\_grid} \\
Werkzeug & Transform, \gui{Grid array} \\
Auswahl & beliebige Objekte \\
Eingabe & 1.~\code{Grid array (3 $\times$ 3): base point (type two numbers to change the
  counts)}\newline
  2.~\code{Grid array (3 $\times$ 3): step in the first direction}\newline
  3.~\code{Grid array (3 $\times$ 3): step in the second direction} \\
Zahlen & vor dem Basispunkt: abwechselnd Anzahl in Richtung~1 und Richtung~2 (ganze Zahlen
  $\geq 1$, einschließlich Original) \\
Bezugspunkt & der Basispunkt \\
Ergebnis & $n_1 \times n_2 - 1$ Kopien im Raster; Kopien übernehmen Name, Layer und Farbe. \\}

\subsection{array\_linear -- Linear array}\label{sec:befehl-array-linear}
\befehlseintrag{%
Name & \cmd{array\_linear}; Alias \code{array} \\
Werkzeug & Transform, \gui{Linear array} \\
Auswahl & beliebige Objekte \\
Eingabe & 1.~\code{Linear array (3 items): base point (type a number to change the count)}\newline
  2.~\code{Linear array (3 items): step to the next item} \\
Zahlen & vor dem Basispunkt: Gesamtzahl (ganze Zahl $\geq 2$, einschließlich Original) \\
Bezugspunkt & der Basispunkt \\
Ergebnis & $n-1$ Kopien, jeweils um den Schritt weiter verschoben. \\}

\subsection{array\_polar -- Polar array}\label{sec:befehl-array-polar}
\befehlseintrag{%
Name & \cmd{array\_polar}; Alias \code{arraypolar} \\
Werkzeug & Transform, \gui{Polar array} \\
Auswahl & beliebige Objekte \\
Eingabe & 1.~\code{Polar array (6 items): centre (type a number to change the count; axis:
  construction-plane normal)}\newline
  2.~\code{Polar array (6 items): total angle in degrees (Enter: 360)} \\
Zahlen & vor dem Zentrum: Gesamtzahl (ganze Zahl $\geq 2$); danach: Gesamtwinkel in Grad
  ($\neq 0$, Betrag $\leq 360$) \\
Bezugspunkt & das Zentrum \\
Ergebnis & $n-1$ Kopien, gleichmäßig über den Gesamtwinkel um die Achse durch das Zentrum
  parallel zur Normalen der Konstruktionsebene verteilt. Beim Vollkreis (\keys{Enter}) fällt
  keine Kopie auf das Original. \\}

\subsection{box -- Box}\label{sec:befehl-box}
\befehlseintrag{%
Name & \cmd{box}; Alias \code{b} \\
Werkzeug & Solids, \gui{Box} \\
Eingabe & 1.~\code{Box: first corner of the base}\newline
  2.~\code{Box: opposite corner of the base}\newline
  3.~\code{Box: height (point or number)} \\
Zahlen & im 3.~Schritt: Höhe mit Vorzeichen (positiv entlang der Normalen der
  Konstruktionsebene) \\
Bezugspunkt & der zuletzt gesetzte Punkt \\
Ergebnis & Quader als geschlossener Körper, Kanten parallel zu den Weltachsen. \\
Hinweise & Die Grundfläche muss in beiden Ebenenrichtungen eine Ausdehnung haben; Höhe 0 wird
  abgelehnt. \\}

\subsection{cap -- Cap planar holes}\label{sec:befehl-cap}
\befehlseintrag{%
Name & \cmd{cap} \\
Werkzeug & Surfaces, \gui{Cap planar holes} \\
Auswahl & Flächen, Körper; sofort \\
Ergebnis & Alle ebenen offenen Randschleifen werden mit ebenen Deckeln geschlossen und alles zu
  einer Hülle vernäht. Das Ergebnis ersetzt das zuerst gewählte Objekt, die übrigen werden
  entfernt. Meldung \code{Capped: closed solid, volume <V>} oder \code{Capped (still open)}. \\
Fehler & \code{Cap planar holes: no planar open loop to cap}, \code{\ldots: there are
  no open edges}, \code{\ldots: select surfaces} \\}

\subsection{circle -- Circle (centre, radius)}\label{sec:befehl-circle}
\befehlseintrag{%
Name & \cmd{circle}; Aliasse \code{c}, \code{ci} \\
Werkzeug & Curves, \gui{Circle (centre, radius)} \\
Eingabe & 1.~\code{Circle: centre}\newline
  2.~\code{Circle: radius (point or number)} \\
Zahlen & im 2.~Schritt: Radius \\
Bezugspunkt & der Mittelpunkt \\
Ergebnis & Vollkreis in der Ebene durch den Mittelpunkt parallel zur Konstruktionsebene. Ein
  Punkt bestimmt den Radius als Abstand in dieser Ebene. \\}

\subsection{circle3p -- Circle (3 points)}\label{sec:befehl-circle3p}
\befehlseintrag{%
Name & \cmd{circle3p}; Alias \code{c3} \\
Werkzeug & Curves, \gui{Circle (3 points)} \\
Eingabe & \code{Circle: point 1 of 3}, \code{Circle: point 2 of 3},
  \code{Circle: point 3 of 3} \\
Bezugspunkt & der zuletzt gesetzte Punkt \\
Ergebnis & Kreis durch die drei Punkte, in deren Ebene. \\
Hinweise & Doppelte oder kollineare Punkte werden abgelehnt. \\}

\subsection{cone -- Cone}\label{sec:befehl-cone}
\befehlseintrag{%
Name & \cmd{cone} \\
Werkzeug & Solids, \gui{Cone} \\
Eingabe & 1.~\code{Cone: centre of the base}\newline
  2.~\code{Cone: radius (point or number)}\newline
  3.~\code{Cone: height (point or number)} \\
Zahlen & Radius; Höhe mit Vorzeichen \\
Bezugspunkt & der Mittelpunkt der Grundfläche \\
Ergebnis & spitzer Kreiskegel, Achse entlang der Normalen der Konstruktionsebene. \\}

\subsection{coons -- Coons patch}\label{sec:befehl-coons}
\befehlseintrag{%
Name & \cmd{coons} \\
Werkzeug & Surfaces, \gui{Coons patch} \\
Auswahl & genau vier Kurven: zwei gegenüberliegende, dann die beiden anderen; sofort \\
Ergebnis & Fläche, die von den vier Kurven begrenzt wird (Laufrichtungen werden angeglichen). \\
Fehler & \code{Coons patch: select four boundary curves: two opposite, then the other two} \\}

\subsection{copy -- Copy}\label{sec:befehl-copy}
\befehlseintrag{%
Name & \cmd{copy}; Aliasse \code{co}, \code{cp} \\
Werkzeug & Transform, \gui{Copy} \\
Auswahl & beliebige Objekte \\
Eingabe & 1.~\code{Copy: point to move from}\newline
  2.~\code{Copy: point to copy to (Enter: done)} -- beliebig oft \\
Bezugspunkt & der erste Punkt \\
Ergebnis & je Zielpunkt eine verschobene Kopie (jeweils ein Rückgängig-Schritt). \\}

\subsection{curvature -- Curvature}\label{sec:befehl-curvature}
\befehlseintrag{%
Name & \cmd{curvature}; Alias \code{curv} \\
Werkzeug & Analyze, \gui{Curvature} \\
Auswahl & genau eine Kurve, Fläche oder ein Körper \\
Eingabe & \code{Curvature: point on the selected curve or surface} \\
Ergebnis & Kurve: \code{Curvature <k>   radius <1/k>} bzw.\ \code{Curvature 0 (straight)}.
  Fläche: \code{k1 <..>   k2 <..>   Gaussian <..>   mean <..>} am nächstgelegenen Punkt; bei
  Körpern auf der nächstgelegenen Teilfläche. \\
Fehler & \code{Curvature: select one curve or surface} \\}

\subsection{cylinder -- Cylinder}\label{sec:befehl-cylinder}
\befehlseintrag{%
Name & \cmd{cylinder}; Alias \code{cyl} \\
Werkzeug & Solids, \gui{Cylinder} \\
Eingabe & 1.~\code{Cylinder: centre of the base}\newline
  2.~\code{Cylinder: radius (point or number)}\newline
  3.~\code{Cylinder: height (point or number)} \\
Zahlen & Radius; Höhe mit Vorzeichen \\
Bezugspunkt & der Mittelpunkt der Grundfläche \\
Ergebnis & Kreiszylinder als geschlossener Körper, Achse entlang der Normalen der
  Konstruktionsebene. \\}

\subsection{difference -- Boolean difference}\label{sec:befehl-difference}
\befehlseintrag{%
Name & \cmd{difference}; Alias \code{bd} \\
Werkzeug & Solids, \gui{Boolean difference} \\
Auswahl & mindestens zwei geschlossene Körper; der \emph{zuerst} gewählte ist der Grundkörper;
  sofort \\
Ergebnis & Grundkörper minus alle anderen. Ersetzt die Eingangskörper; übernimmt Name, Layer,
  Farbe des Grundkörpers. \\
Fehler & \code{Select the solid to cut first, then the solids to subtract};
  \code{Boolean difference failed: Booleans need closed solids};
  \code{\ldots: the result is empty} \\}

\subsection{distance -- Distance}\label{sec:befehl-distance}
\befehlseintrag{%
Name & \cmd{distance}; Aliasse \code{dist}, \code{di} \\
Werkzeug & Analyze, \gui{Distance} \\
Eingabe & 1.~\code{Distance: first point}\newline 2.~\code{Distance: second point} \\
Bezugspunkt & der erste Punkt \\
Ergebnis & Meldung \code{Distance <d>   (dx <..>, dy <..>, dz <..>)}; bleibt in der
  Statusleiste stehen. \\}

\subsection{explode -- Explode}\label{sec:befehl-explode}
\befehlseintrag{%
Name & \cmd{explode}; Alias \code{x} \\
Werkzeug & Edit, \gui{Explode} \\
Auswahl & Körper/Hüllen, Kurven; sofort \\
Ergebnis & Hüllen zerfallen in ihre Flächen, Kurven an ihren Knicken in Teilkurven
  (Polylinien in Strecken; Kreise und Bögen über \SI{90}{\degree} in ihre Teilbögen). Meldung
  mit der Anzahl der Teile. \\
Fehler & \code{Explode: nothing to explode} \\}

\subsection{extrude -- Extrude}\label{sec:befehl-extrude}
\befehlseintrag{%
Name & \cmd{extrude}; Alias \code{ext} \\
Werkzeug & Surfaces, \gui{Extrude} \\
Auswahl & eine oder mehrere Kurven \\
Eingabe & 1.~\code{Extrude: distance along the construction-plane normal, or start of the
  direction}\newline
  2.~(nach einem Punkt) \code{Extrude: end of the direction} \\
Zahlen & im 1.~Schritt: Distanz ($\neq 0$) entlang der Normalen der Konstruktionsebene der
  Ansicht, in der Sie arbeiten (aktive Ansicht beim Start, danach die Ansicht unter dem
  Cursor): Top und Perspective $+z$, Front $-y$, Left $-x$; negativ in Gegenrichtung. Im
  2.~Schritt: Länge in Richtung des Cursors \\
Bezugspunkt & der Anfangspunkt der Richtung \\
Ergebnis & je Kurve eine Extrusionsfläche (offen; mit \cmd{cap} schließen). \\
Fehler & \code{Extrude: select one or more curves}; \code{Extrude: the distance must not be 0} \\}

\subsection{fillet\_curve -- Fillet curves}\label{sec:befehl-fillet-curve}
\befehlseintrag{%
Name & \cmd{fillet\_curve}; Alias \code{fillet} \\
Werkzeug & Curves, \gui{Fillet curves} \\
Auswahl & eine oder zwei ebene Kurven \\
Eingabe & \code{Fillet: radius (two curves: fillet between them; one curve: all its corners)}
  -- nur Zahl \\
Ergebnis & Zwei Kurven: neuer Verrundungsbogen, beide Kurven gekürzt. Eine Kurve: alle Ecken
  verrundet, die Kurve wird ersetzt. \\
Fehler & \code{Fillet curves: select one or two planar curves}; \code{Fillet: radius <r> too
  large for these curves}; \code{Fillet: the radius must be positive}; \code{Fillet: the
  curves do not meet – extend them or pick two curves that cross}; \code{Fillet: the curves
  are tangent where they meet (nothing to round)}; \code{Fillet: the curve is not planar};
  \code{Fillet: the curve has no corners}. Der Befehl wartet danach auf einen anderen
  Radius. \\}

\subsection{fillet\_surface -- Fillet surfaces}\label{sec:befehl-fillet-surface}
\befehlseintrag{%
Name & \cmd{fillet\_surface}; Alias \code{filletsrf} \\
Werkzeug & Surfaces, \gui{Fillet surfaces} \\
Auswahl & genau zwei unbeschnittene Flächen \\
Eingabe & \code{Fillet surfaces: radius} -- nur Zahl ($\neq 0$; das Vorzeichen wird
  ignoriert) \\
Ergebnis & Rundungsfläche zwischen beiden Flächen; die Seite wird automatisch bestimmt. Die
  Flächen selbst bleiben unverändert. \\
Fehler & \code{Fillet surfaces: select two (untrimmed) surfaces} \\}

\subsection{gordon -- Gordon surface}\label{sec:befehl-gordon}
\befehlseintrag{%
Name & \cmd{gordon}; Alias \code{network} \\
Werkzeug & Surfaces, \gui{Gordon surface} \\
Auswahl & ein Kurvennetz aus mindestens $2 \times 2$ sich kreuzenden Kurven, beliebige
  Reihenfolge; sofort \\
Ergebnis & Fläche durch das Netz. Die Scharen werden automatisch gebildet (Kurven, die die
  zuerst gewählte schneiden, und die übrigen) und sortiert. \\
Fehler & \code{Gordon surface: select a curve network (at least 2 $\times$ 2 crossing
  curves)}; \code{Gordon surface: the curves do not form a network} \\}

\subsection{helix -- Helix}\label{sec:befehl-helix}
\befehlseintrag{%
Name & \cmd{helix}; Aliasse \code{h}, \code{he} \\
Werkzeug & Curves, \gui{Helix} \\
Eingabe & 1.~\code{Helix (5 turns): axis start (type a number to change the turns)}\newline
  2.~\code{Helix (5 turns): axis end}\newline
  3.~\code{Helix (5 turns): radius (point or number)} \\
Zahlen & in Schritt 1 und 2: Windungszahl ($>0$); in Schritt 3: Radius \\
Bezugspunkt & der zuletzt gesetzte Punkt \\
Ergebnis & Schraubenlinie; Höhe = Achslänge, Steigung = Höhe / Windungen. Ein Radiuspunkt
  legt auch den Startpunkt fest. \\
Hinweise & In einer Ansicht wird \keys{H} nur dann in die Befehlszeile geschrieben, wenn
  nichts ausgewählt ist; mit Auswahl blendet \keys{H} die Auswahl aus (Hide). \\}

\subsection{intersection -- Boolean intersection}\label{sec:befehl-intersection}
\befehlseintrag{%
Name & \cmd{intersection}; Alias \code{bi} \\
Werkzeug & Solids, \gui{Boolean intersection} \\
Auswahl & mindestens zwei geschlossene Körper; sofort \\
Ergebnis & gemeinsamer Bereich aller Körper; ersetzt die Eingangskörper. \\
Fehler & \code{Select at least two solids}; \code{Boolean intersection failed: \ldots} \\
Hinweise & Nicht verwechseln mit \cmd{intersection\_check} (Schnittkurven). \\}

\subsection{intersection\_check -- Intersect}\label{sec:befehl-intersection-check}
\befehlseintrag{%
Name & \cmd{intersection\_check}; Alias \code{intersect} \\
Werkzeug & Analyze, \gui{Intersection check} \\
Auswahl & mindestens zwei Kurven/Flächen/Körper; sofort \\
Ergebnis & Schnittpunkte (Kurve--Kurve, Kurve--Fläche) als Punktobjekte, Schnittkurven
  (Fläche--Fläche) als Kurven. Meldung \code{Intersections: <n> curve(s), <m> point(s)}.
  Flächen desselben Körpers werden nicht geschnitten. \\
Fehler & \code{Intersect: select at least two objects} \\}

\subsection{join -- Join}\label{sec:befehl-join}
\befehlseintrag{%
Name & \cmd{join}; Alias \code{j} \\
Werkzeug & Edit, \gui{Join} \\
Auswahl & mindestens zwei Kurven \emph{oder} mindestens zwei Flächen/Körper; sofort \\
Ergebnis & Kurven: eine verkettete Kurve (beliebige Reihenfolge und Richtung). Flächen: eine
  vernähte Hülle, Meldung \code{Joined <n> faces, closed solid} bzw.\ \code{\ldots, open}.
  Ersetzt das zuerst gewählte Objekt. \\
Fehler & \code{Join: select either curves or surfaces}; \code{\ldots: select at least
  two objects}; \code{\ldots: the curves do not form one chain} \\}

\subsection{line -- Line}\label{sec:befehl-line}
\befehlseintrag{%
Name & \cmd{line}; Aliasse \code{l}, \code{li} \\
Werkzeug & Draw, \gui{Line} \\
Eingabe & 1.~\code{Line: start point}\newline 2.~\code{Line: end point (or length)} \\
Zahlen & im 2.~Schritt: Länge in Richtung des Cursors \\
Bezugspunkt & der Startpunkt \\
Ergebnis & eine Strecke. \\}

\subsection{loft -- Loft}\label{sec:befehl-loft}
\befehlseintrag{%
Name & \cmd{loft} \\
Werkzeug & Surfaces, \gui{Loft} \\
Auswahl & zwei oder mehr Querschnittskurven in Auswahlreihenfolge; sofort \\
Ergebnis & glatte Fläche durch alle Querschnitte. \\
Fehler & \code{Loft: select two or more section curves (in order)}; \code{Loft: <Grund>} \\}

\subsection{mass\_properties -- Volume / mass properties}\label{sec:befehl-mass-properties}
\befehlseintrag{%
Name & \cmd{mass\_properties}; Alias \code{volume} \\
Werkzeug & Analyze, \gui{Volume / mass properties} \\
Auswahl & geschlossene Körper (andere Objekte werden übergangen); sofort \\
Ergebnis & Meldung \code{Volume <V>   centroid (<x>, <y>, <z>)}: Gesamtvolumen und
  gemeinsamer Schwerpunkt. \\
Fehler & \code{Volume / mass properties: select closed solids} \\}

\subsection{mirror -- Mirror}\label{sec:befehl-mirror}
\befehlseintrag{%
Name & \cmd{mirror}; Alias \code{mi} \\
Werkzeug & Transform, \gui{Mirror} \\
Auswahl & beliebige Objekte \\
Eingabe & 1.~\code{Mirror: start of the mirror line (copies, plane through the line and the
  construction-plane normal)}\newline
  2.~\code{Mirror: end of the mirror line} \\
Bezugspunkt & der Anfang der Spiegellinie \\
Ergebnis & gespiegelte Kopien; Spiegelebene durch die Linie und die Normale der
  Konstruktionsebene. \\}

\subsection{move -- Move}\label{sec:befehl-move}
\befehlseintrag{%
Name & \cmd{move}; Alias \code{m} \\
Werkzeug & Transform, \gui{Move} \\
Auswahl & beliebige Objekte \\
Eingabe & 1.~\code{Move: point to move from}\newline 2.~\code{Move: point to move to} \\
Bezugspunkt & der erste Punkt \\
Ergebnis & die Objekte selbst werden verschoben. \\}

\subsection{offset\_curve -- Offset curve}\label{sec:befehl-offset-curve}
\befehlseintrag{%
Name & \cmd{offset\_curve}; Alias \code{offset} \\
Werkzeug & Curves, \gui{Offset curve} \\
Auswahl & eine oder mehrere Kurven \\
Eingabe & \code{Offset curve (distance 1): click the side (type a number to change the
  distance)} \\
Zahlen & Abstand ($>0$), beliebig oft vor dem Klick \\
Ergebnis & je Kurve eine Parallelkurve auf der angeklickten Seite, in der Konstruktionsebene
  der Ansicht, in der geklickt wurde. \\
Fehler & \code{Offset curve: select one or more curves}; abgelehnte Eingaben (der Befehl
  wartet weiter): \code{Offset curve: the distance must be positive (click the side)};
  \code{\ldots: the curve does not lie in the construction plane of this view – click the side
  in a view parallel to the curve}; \code{\ldots: the curve has corners – Explode it and
  offset the pieces, or fillet the corners first}; \code{\ldots: distance <d> is larger than
  the radius of curvature on that side} \\}

\subsection{offset\_surface -- Offset surface}\label{sec:befehl-offset-surface}
\befehlseintrag{%
Name & \cmd{offset\_surface}; Alias \code{offsetsrf} \\
Werkzeug & Surfaces, \gui{Offset surface} \\
Auswahl & eine oder mehrere unbeschnittene Flächen \\
Eingabe & \code{Offset surface: distance along the normal (negative: other side)} -- nur
  Zahl ($\neq 0$) \\
Ergebnis & je Fläche eine Parallelfläche. \\
Fehler & \code{Offset surface: select (untrimmed) surfaces} \\}

\subsection{plane -- Plane}\label{sec:befehl-plane}
\befehlseintrag{%
Name & \cmd{plane} \\
Werkzeug & Surfaces, \gui{Plane} \\
Eingabe & 1.~\code{Plane: first corner}\newline 2.~\code{Plane: opposite corner} \\
Bezugspunkt & die erste Ecke \\
Ergebnis & rechteckige ebene Fläche, Kanten parallel zu den Achsen der Konstruktionsebene. \\}

\subsection{point -- Point}\label{sec:befehl-point}
\befehlseintrag{%
Name & \cmd{point}; Aliasse \code{pt}, \code{po} \\
Werkzeug & Draw, \gui{Point} \\
Eingabe & \code{Point: location (Enter / right-click: done)} -- beliebig oft \\
Bezugspunkt & der zuletzt gesetzte Punkt \\
Ergebnis & je Eingabe ein Punktobjekt (jeweils ein Rückgängig-Schritt). \\}

\subsection{polygon -- Polygon}\label{sec:befehl-polygon}
\befehlseintrag{%
Name & \cmd{polygon}; Alias \code{pol} \\
Werkzeug & Draw, \gui{Polygon} \\
Eingabe & 1.~\code{Polygon (6 sides): centre (type a number to change the sides)}\newline
  2.~\code{Polygon (6 sides): corner point or radius} \\
Zahlen & vor dem Mittelpunkt: Eckenzahl (ganze Zahl 3 bis 1000); danach: Umkreisradius (erste
  Ecke in $u$-Richtung) \\
Bezugspunkt & der Mittelpunkt \\
Ergebnis & regelmäßiges geschlossenes Vieleck in der Konstruktionsebene. \\}

\subsection{polyline -- Polyline}\label{sec:befehl-polyline}
\befehlseintrag{%
Name & \cmd{polyline}; Alias \code{pl} \\
Werkzeug & Draw, \gui{Polyline} \\
Eingabe & 1.~\code{Polyline: start point}\newline
  2.~\code{Polyline: next point (Enter: done, snap to the first point: close)} -- beliebig
  oft \\
Bezugspunkt & der letzte Punkt \\
Ergebnis & \keys{Enter}/Rechtsklick: offene Polylinie (ab 2 Punkten). Erneut der Startpunkt (ab
  3 vorhandenen Punkten): geschlossene Polylinie. Die gesetzten Punkte sind während des
  Zeichnens als \gui{End} fangbar, so dass sich der Startpunkt exakt treffen lässt. \\}

\subsection{rectangle -- Rectangle}\label{sec:befehl-rectangle}
\befehlseintrag{%
Name & \cmd{rectangle}; Aliasse \code{rec}, \code{rect} \\
Werkzeug & Draw, \gui{Rectangle} \\
Eingabe & 1.~\code{Rectangle: first corner}\newline 2.~\code{Rectangle: opposite corner} \\
Bezugspunkt & die erste Ecke \\
Ergebnis & geschlossenes Rechteck (eine Kurve), Kanten parallel zu den Achsen der
  Konstruktionsebene. \\}

\subsection{revolve -- Revolve}\label{sec:befehl-revolve}
\befehlseintrag{%
Name & \cmd{revolve}; Alias \code{rev} \\
Werkzeug & Surfaces, \gui{Revolve} \\
Auswahl & eine oder mehrere Kurven (Profil) \\
Eingabe & 1.~\code{Revolve: start of the axis}\newline
  2.~\code{Revolve: end of the axis}\newline
  3.~\code{Revolve: angle in degrees (Enter: 360)} \\
Zahlen & im 3.~Schritt: Winkel ($\neq 0$, Betrag $\leq 360$); positiv gegen den
  Uhrzeigersinn um die Achsrichtung \\
Bezugspunkt & der zuletzt gesetzte Achspunkt \\
Ergebnis & je Profilkurve eine Rotationsfläche. \\
Fehler & \code{Revolve: select one or more curves}; \code{Revolve: the end of the axis must
  differ from its start}; \code{Revolve: the angle must be between -360 and 360 degrees and
  not 0}; \code{Revolve: type the angle in degrees, or Enter for 360} (ein Punkt statt des
  Winkels) \\}

\subsection{rotate -- Rotate}\label{sec:befehl-rotate}
\befehlseintrag{%
Name & \cmd{rotate}; Alias \code{ro} \\
Werkzeug & Transform, \gui{Rotate} \\
Auswahl & beliebige Objekte \\
Eingabe & 1.~\code{Rotate: centre (axis: normal of the construction plane)}\newline
  2.~\code{Rotate: angle in degrees or first reference point}\newline
  3.~(nach einem Bezugspunkt) \code{Rotate: second reference point} \\
Zahlen & im 2.~Schritt: Winkel in Grad \\
Bezugspunkt & das Zentrum \\
Ergebnis & die Objekte werden um die Achse durch das Zentrum parallel zur Normalen der
  Konstruktionsebene gedreht. \\}

\subsection{scale -- Scale}\label{sec:befehl-scale}
\befehlseintrag{%
Name & \cmd{scale}; Alias \code{sc} \\
Werkzeug & Transform, \gui{Scale} \\
Auswahl & beliebige Objekte \\
Eingabe & 1.~\code{Scale: base point}\newline
  2.~\code{Scale: factor or first reference point}\newline
  3.~(nach einem Bezugspunkt) \code{Scale: second reference point} \\
Zahlen & im 2.~Schritt: Faktor ($>0$) \\
Bezugspunkt & der Basispunkt \\
Ergebnis & gleichmäßige Skalierung um den Basispunkt; mit Bezugspunkten ist der Faktor das
  Abstandsverhältnis (Punkt~2 zu Punkt~1). \\}

\subsection{sphere -- Sphere}\label{sec:befehl-sphere}
\befehlseintrag{%
Name & \cmd{sphere}; Alias \code{sph} \\
Werkzeug & Solids, \gui{Sphere} \\
Eingabe & 1.~\code{Sphere: centre}\newline 2.~\code{Sphere: radius (point or number)} \\
Bezugspunkt & der Mittelpunkt \\
Ergebnis & Kugel als geschlossener Körper. \\}

\subsection{spline -- Interpolated spline}\label{sec:befehl-spline}
\befehlseintrag{%
Name & \cmd{spline}; Alias \code{sp} \\
Werkzeug & Curves, \gui{Interpolated spline} \\
Eingabe & 1.~\code{Spline: first point}\newline
  2.~\code{Spline: next point (Enter: done)} -- beliebig oft \\
Bezugspunkt & der letzte Punkt \\
Ergebnis & mit \keys{Enter}: kubische Spline-Kurve durch alle Punkte (ab 2 Punkten; bei 2 bzw.\
  3 Punkten Grad 1 bzw.\ 2). \\}

\subsection{split -- Split}\label{sec:befehl-split}
\befehlseintrag{%
Name & \cmd{split} \\
Werkzeug & Edit, \gui{Split} \\
Auswahl & (a) genau eine Kurve, oder (b) zuerst eine Fläche, dann eine Schnittfläche \\
Eingabe & (a) \code{Split: click the point on the curve}; (b) sofort \\
Ergebnis & (a) zwei Kurven; (b) die Teilflächen. Das Original wird ersetzt. \\
Fehler & \code{Split: select one curve, or a surface and then a cutting surface};
  \code{Split: the surfaces do not intersect}; Punkt am Kurvenende: \code{Split: the point is
  at an end of the curve} (der Befehl wartet auf einen anderen Punkt) \\}

\subsection{sweep1 -- Sweep 1 rail}\label{sec:befehl-sweep1}
\befehlseintrag{%
Name & \cmd{sweep1} \\
Werkzeug & Surfaces, \gui{Sweep 1 rail} \\
Auswahl & genau zwei Kurven: Profil, dann Leitkurve; sofort \\
Ergebnis & Fläche: Profil entlang der Leitkurve geführt. \\
Fehler & \code{Sweep 1 rail: select the profile, then the rail} \\}

\subsection{sweep2 -- Sweep 2 rails}\label{sec:befehl-sweep2}
\befehlseintrag{%
Name & \cmd{sweep2} \\
Werkzeug & Surfaces, \gui{Sweep 2 rails} \\
Auswahl & genau drei Kurven: Profil, Leitkurve~1, Leitkurve~2; sofort \\
Ergebnis & Fläche: Profil entlang zweier Leitkurven geführt. \\
Fehler & \code{Sweep 2 rails: select the profile, then rail 1 and rail 2} \\}

\subsection{torus -- Torus}\label{sec:befehl-torus}
\befehlseintrag{%
Name & \cmd{torus}; Alias \code{tor} \\
Werkzeug & Solids, \gui{Torus} \\
Eingabe & 1.~\code{Torus: centre}\newline
  2.~\code{Torus: major radius (point or number)}\newline
  3.~\code{Torus: tube radius (point or number)} \\
Bezugspunkt & der Mittelpunkt \\
Ergebnis & Torus mit Mittelebene parallel zur Konstruktionsebene. \\
Hinweise & Der Rohrradius muss kleiner als der große Radius sein; ein Punkt bestimmt ihn als
  Abstand vom Mittelkreis. \\}

\subsection{trim -- Trim}\label{sec:befehl-trim}
\befehlseintrag{%
Name & \cmd{trim}; Alias \code{tr} \\
Werkzeug & Edit, \gui{Trim} \\
Auswahl & zuerst das zu stutzende Objekt (Kurve oder einzelne Fläche), dann die Schneidobjekte
  (Kurven/Flächen bzw.\ eine Schnittfläche) \\
Eingabe & \code{Trim: click the part of the first selected object to remove} \\
Ergebnis & Das Objekt wird an den Schnitten geteilt, das angeklickte Stück entfernt, die übrigen
  bleiben. \\
Fehler & \code{Trim: select the object to trim first, then the cutting objects}; kein Schnitt:
  \code{Trim: the cutting objects do not cross the curve} bzw.\ \code{Trim: the surfaces do not
  intersect}; Klick neben die Teile einer Fläche: \code{Trim: click on the part to remove} (der
  Befehl wartet jeweils auf einen neuen Klick) \\}

\subsection{union -- Boolean union}\label{sec:befehl-union}
\befehlseintrag{%
Name & \cmd{union}; Alias \code{bu} \\
Werkzeug & Solids, \gui{Boolean union} \\
Auswahl & mindestens zwei geschlossene Körper; sofort \\
Ergebnis & Vereinigung aller Körper; ersetzt die Eingangskörper, übernimmt Name, Layer und
  Farbe des zuerst gewählten. \\
Fehler & \code{Select at least two solids}; \code{Boolean union failed: \ldots} \\}

\fussy
